% !PRINTED 733
\weberpage{733}
Setzt man
\[
n\varkappa\Sigma\delta_s|g_s| = g,
\]
so ist um so mehr, da die Summe der absoluten Werte nicht kleiner ist, als der absolute Wert der Summe:

\weq{14}{-g + \alpha u < \Sigma g_s\lambda_s < g + \alpha u.}
Man setze nun
\[
-g + \alpha u = a, \qquad g + \alpha u = a'
\]
\[
\alpha\delta = a' - a = 2g
\]
und erhält:

\weq{15}{a < \Sigma g_s\lambda_s < a'.}
Ersetzt man $u$ durch $u' = u + \delta$, so geht $a$ in $a'$ über (das wäre bei der Formulierung des Textes nicht der Fall), und wenn man $a'' = a' + \alpha\delta = g + \alpha u'$ setzt, und wenn $\lambda_s'$ ebenso von $u'$ abhängt wie $\lambda_s$ von $u$, so ergibt sich
\[
a' < \Sigma g_s\lambda_s' < a''.
\]
Durch Fortsetzung dieses Verfahrens wird eine Reihe von ganzen Zahlen
\[
\eta, \eta', \eta'', \eta''', \ldots
\]
den folgenden Bedingungen gemäß bestimmt.

Von hier kann wie im Text Seite 698 oben fortgefahren werden. Der Text auf S. 697 wird richtig durch folgende vier Korrekturen:

Seite 697, Zeile 13 v. u. lies $nk\Sigma\delta_s|g_s|$ statt $k\Sigma\delta_s\gamma_s$.

\glqq{} 11 v. u. \glqq{} $-g$ statt $-ng$.

\glqq{} 3 v. u. \glqq{} $-g$ \glqq{} $-ng$.

\glqq{} 1 v. u. \glqq{} $2g$ \glqq{} $(n+1)g$.

\begin{center}
\rule{3cm}{0.4pt}
\end{center}

\begin{center}
\textbf{Berichtigungen zum dritten Bande.}
\end{center}

\begin{center}
\rule{3cm}{0.4pt}
\end{center}

Seite 28, Formel (23) lies $K^2$ statt $K_2$.

\glqq{} 124, Formelsystem (8) lies $u$ statt $u_1$.

\glqq{} 180, Zeile 9 v. u. lies $\alpha^2\beta\delta$ \glqq{} $\alpha^2\beta\gamma$.

\glqq{} 185, letzte Zeile und Seite 186, Formel (6), (7) lies $g_2^2$ statt $g_3^2$.

\glqq{} 333, Zeile 8 v. o. lies äquivalent, wenn statt äquivalent. Wenn.

\glqq{} 340, Zeile 4 v. u. lies $a_{22}$ statt $a_{12}$.

\glqq{} 394, \glqq{} 18 v. o. b) ist $-5$ zu streichen.

\glqq{} 525, Formel (1) lies $Ax^2$ statt $Ax$.

\begin{center}
\makebox[3cm]{\dotfill}
\end{center}
